\documentclass[11pt,a4paper]{article}
\usepackage[margin=27mm]{geometry}
\usepackage{amsmath,amssymb,amsthm,fontspec,unicode-math,hyperref,booktabs,longtable}
\usepackage{array,calc,fvextra}
\setmainfont{STIXGeneral}
\setmathfont{STIX Two Math}
\setmonofont{Menlo}
\hypersetup{hidelinks}
\setlength{\emergencystretch}{2em}
\providecommand{\tightlist}{\setlength{\itemsep}{0pt}\setlength{\parskip}{0pt}}
\newtheorem{theorem}{Theorem}
\newtheorem{lemma}[theorem]{Lemma}
\newtheorem{corollary}[theorem]{Corollary}
\newcommand{\gen}[1]{\langle #1\rangle}
\newcommand{\Sym}{\operatorname{Sym}}
\newcommand{\Aut}{\operatorname{Aut}}
\newcommand{\Fix}{\operatorname{Fix}}

\title{Odd kernels and cyclotomic modules}
\author{The human\_compression project}
\date{Canonical Topic F: common theory}
\begin{document}
\maketitle
\begin{abstract}
An involution outside an odd normal subgroup is controlled by its orbit
under a Klein four complement. In the abelian case, this orbit has two
cyclic character components. For nonabelian constructions, quotients of
the cyclotomic lattice give explicit finite coordinates, automorphisms
and presentation relations. We give the complete elementary arguments
and the subgroup intersection tests needed to turn these constructions
into regular polyhedra.
\end{abstract}
\section{Marked groups and intersections}
% BEGIN UF.Conventions
We use $x^g=g^{-1}xg$ and $[x,y]=x^{-1}y^{-1}xy$.
A rank-three string marking of a finite group $G$ is a generating triple
of distinct nonidentity involutions $r_0,r_1,r_2$ whose endpoints commute.
Write $G_I=\gen{r_i:i\in I}$ for $I\subseteq\{0,1,2\}$.
The marking is a \emph{string C-group} if
\[
 G_I\cap G_J=G_{I\cap J}\qquad(I,J\subseteq\{0,1,2\}).
\]
Its type is $\{k_1,k_2\}$, where $k_1=|r_0r_1|$ and
$k_2=|r_1r_2|$. We call it nondegenerate when $k_1,k_2\geq3$,
and tight when $|G|=2k_1k_2$. A marked isomorphism preserves the
three generators; duality reverses them. We retain the classical
correspondence between string C-groups and abstract regular polyhedra
when using geometric terminology.

% END UF.Conventions

% BEGIN UF.RankThree
\begin{lemma}[Rank-three intersection test]\label{UF.RankThree}
A nondegenerate rank-three string marking is a string C-group if and
only if
\begin{equation}\label{eq:rankthree}
 \gen{r_0,r_1}\cap\gen{r_1,r_2}=\gen{r_1}.
\end{equation}
In that case $2k_1k_2\leq |G|$.
\end{lemma}
\begin{proof}
Necessity is part of the intersection property. For sufficiency set
$H_{ij}=\gen{r_i,r_j}$. The endpoint group consists of
$1,r_0,r_2,r_0r_2$. If its intersection with $H_{01}$ contains $r_2$
or $r_0r_2$, then $r_2\in H_{01}\cap H_{12}=\gen{r_1}$,
contrary to distinctness. Hence
$H_{02}\cap H_{01}=\gen{r_0}$; dually
$H_{02}\cap H_{12}=\gen{r_2}$. These equalities and
\eqref{eq:rankthree} also exclude a generator from the group on the
other two. Singleton intersections follow, and pairs involving the
empty or full index set are immediate. This accounts for every pair
of subsets. Finally $|H_{01}|=2k_1$ and $|H_{12}|=2k_2$, so
\[
 |H_{01}H_{12}|=\frac{|H_{01}|\,|H_{12}|}{|H_{01}\cap H_{12}|}
 =2k_1k_2\leq |G|.
\]
\end{proof}
% END UF.RankThree

% BEGIN UF.PrimeIndex
\begin{lemma}[Prime-index intersection]\label{UF.PrimeIndex}
Let $C,H,J$ be subgroups of a group, with $C\leq H\cap J$ and
$[H:C]=p$ prime. If $H\not\leq J$, then $H\cap J=C$.
In particular, this holds if $\gen{H,J}=G$ and $J<G$.
\end{lemma}
\begin{proof}
The index tower gives
$p=[H:H\cap J][H\cap J:C]$. The first factor is not one,
so the second is one. In the final situation $H\leq J$ would give
$G=\gen{H,J}=J$, contrary to the hypothesis.
\end{proof}
For later products it is useful to record the elementary identities
\begin{equation}\label{eq:reflections}
 (sh)^2=1,\qquad (sh)(bh)=sb^{-1},
\end{equation}
whenever $h^2=1$, $hsh=s^{-1}$ and $hbh=b^{-1}$.
The first equality uses only the condition involving $s$; the second
uses only the condition involving $b$.
% END UF.PrimeIndex

\section{An inverted orbit and its characters}
% BEGIN UF.Orbit
\begin{lemma}[One inverted orbit]\label{UF.Orbit}
Suppose $G=P\rtimes K$, where $P$ has odd order and
$K=\gen{a,b}$ is a Klein four group. Suppose $G=\gen{K,r}$,
where $r$ is a nonidentity involution. Write $r=ut$ with
$u\in P$, $t\in K$. Then $t\ne1$, $u^t=u^{-1}$, and
\[
 P=\gen{u,u^k}\qquad\text{for every }k\in K\setminus\gen{t}.
\]
\end{lemma}
\begin{proof}
An odd-order group has no nonidentity involution, so $t\ne1$.
The equation $(ut)^2=1$ gives $u^t=u^{-1}$. Put
$L=\gen{u,u^k}$. Since $t$ commutes with $k$, it inverts both
displayed generators of $L$. The element $k$ exchanges them. Thus
$L$ is $K$-invariant and $LK$ is a subgroup containing $K$ and $r$.
It equals $G$, and intersecting with $P$ gives $L=P$.
\end{proof}
This proof says that $t$ inverts the two indicated generators. It does
not assert that inversion is an automorphism of a nonabelian group.
% END UF.Orbit

% BEGIN UF.Characters
\begin{lemma}[Two cyclic characters]\label{UF.Characters}
In Lemma~\ref{UF.Orbit}, suppose $P$ is abelian and noncyclic.
Then it is the direct product of two nontrivial cyclic subgroups
$P_\chi$, indexed by the two characters
$\chi:K\longrightarrow\{1,-1\}$ with $\chi(t)=-1$.
The projection of $u$ generates each $P_\chi$.
\end{lemma}
\begin{proof}
Write $P$ additively. Multiplication by four is an automorphism of $P$.
The four endomorphisms
\[
 e_{\epsilon,\delta}=\tfrac14(1+\epsilon a)(1+\delta b),
 \qquad\epsilon,\delta\in\{1,-1\},
\]
are orthogonal idempotents with sum the identity: expand the products
using $a^2=b^2=1$ and $ab=ba$. Their images give the four character
components, on which $a,b$ act by $\epsilon,\delta$ respectively.
The element $t$ inverts $u$ and commutes with $K$, so it inverts the
whole $K$-orbit of $u$. That orbit generates $P$ by
Lemma~\ref{UF.Orbit}; hence $t$ acts as $-1$ on $P$.
Every component with $\chi(t)=1$ is therefore annihilated by two and
is zero. For either remaining character, the projection of each orbit
element is plus or minus $e_\chi u$. Its component is consequently
cyclic. If one component were zero then $P$ would be cyclic, so both
are nontrivial.
\end{proof}
% END UF.Characters

\section{The complete cyclotomic lattice}
Let $p$ be an odd prime, let $\zeta$ be a primitive $p$th root of unity,
and put $O=\mathbb Z[\zeta]$ and $\pi=\zeta-1$.
% BEGIN UF.Lattice
\begin{theorem}[Cyclotomic coordinates]\label{UF.Lattice}
For every $n\geq1$, the ring $A_n=O/\pi^nO$ has $p^n$ elements.
Write uniquely
\[
 n=(e-1)(p-1)+r,\qquad e\geq1,\quad 1\leq r\leq p-1.
\]
As an additive group with the indicated ordered generators,
\begin{equation}\label{eq:coordinates}
 A_n=\gen{1}\oplus\gen{\pi}\oplus\cdots\oplus\gen{\pi^{p-2}}
 \cong (\mathbb Z/p^e\mathbb Z)^r
 \oplus(\mathbb Z/p^{e-1}\mathbb Z)^{p-1-r}.
\end{equation}
A modulus-one factor is trivial. Complex conjugation induces an
involutory automorphism of $A_n$. Multiplication by $\zeta$ has order
$p$ when $n\geq2$, and for $j\not\equiv0\pmod p$ it satisfies
$\sum_{i=0}^{p-1}\zeta^{ij}=0$.
\end{theorem}
\begin{proof}
The polynomial
\[
 \Phi_p(1+X)=p+\binom p2X+\cdots+
 \binom p{p-1}X^{p-2}+X^{p-1}
\]
is Eisenstein at $p$. Thus $O$ is an integral domain with
$\mathbb Z$-basis $1,\pi,\ldots,\pi^{p-2}$. Substitution
$\zeta=1$ gives $O/\pi O\cong\mathbb F_p$.
Multiplication by $\pi^j$ induces an isomorphism
$O/\pi O\longrightarrow\pi^jO/\pi^{j+1}O$:
surjectivity is immediate and injectivity follows by cancellation
in the domain. The successive indices are all $p$, whence
\begin{equation}\label{eq:latticeorder}
 |O/\pi^nO|=p^n.
\end{equation}

The identity $p=\prod_{j=1}^{p-1}(1-\zeta^j)$ gives
$pO=\pi^{p-1}O$. Here equality, rather than inclusion alone, follows
because each ratio $(1-\zeta^j)/(1-\zeta)$ is a unit of $O$.
It belongs to $O$ by a geometric sum. The automorphism
$\zeta\mapsto\zeta^j$ identifies the norms of numerator and
denominator, both equal to $p$, so the ratio has norm one.
Its multiplication matrix on the integral basis has determinant one;
the adjugate gives an integral inverse and hence a unit of $O$.

For $1\leq r\leq p-1$ define the integral lattice
\[
 L_r=p\mathbb Z\,1+\cdots+p\mathbb Z\,\pi^{r-1}
       +\mathbb Z\,\pi^r+\cdots+\mathbb Z\,\pi^{p-2}.
\]
Every displayed generator is in $\pi^rO$, using
$pO=\pi^{p-1}O$ for the first group of generators.
Both $L_r$ and $\pi^rO$ have index $p^r$ in $O$, the latter by
\eqref{eq:latticeorder}; thus $L_r=\pi^rO$.
Now $\pi^nO=p^{e-1}\pi^rO$, which gives every coordinate modulus
in \eqref{eq:coordinates}. In particular it gives the individual
orders, including trivial listed generators when $e=1$.

The equation $\overline\pi=-\zeta^{-1}\pi$ shows that conjugation
preserves $\pi^nO$. Multiplication by $\zeta$ has $p$th power the
identity. For $n\geq2$ it is not the identity, since its difference
from the identity sends $1$ to $\pi\notin\pi^nO$; the latter follows
from the strict successive indices above. Its order is therefore $p$.
Finally $1+\zeta+\cdots+\zeta^{p-1}=0$; multiplication of exponents
by a nonzero residue $j$ permutes these terms.
\end{proof}
% END UF.Lattice

% BEGIN UF.Conjugation
\begin{corollary}[Conjugation and displacement]\label{UF.Conjugation}
Let $A_n$ be written additively and let an element $\beta$ of order
$p$ act by $v^\beta=\zeta v$. Put $s_i=\pi^{i-1}$ for $i\geq1$.
Then $[s_i,\beta]=s_{i+1}$ in $A_n\rtimes\gen{\beta}$, and
\begin{align}
 \overline{s_i}&=(-1)^{i-1}\zeta^{1-i}s_i,
       &&i\geq1,\label{eq:conjugation}\\
 \zeta^{1-i}s_i&=\sum_{j=i}^{p+1}
       \binom{p+1-i}{j-i}s_j,
       &&1\leq i\leq p.\label{eq:binomial}
\end{align}
The first expression on the left of the commutator identity is group
notation; the two displayed equations are additive ring identities.
\end{corollary}
\begin{proof}
The commutator is $(\zeta-1)s_i=\pi s_i$.
Equation~\eqref{eq:conjugation} follows by raising
$\overline\pi=-\zeta^{-1}\pi$ to the power $i-1$.
For \eqref{eq:binomial}, use $\zeta^{1-i}=\zeta^{p+1-i}$ and
expand $(1+\pi)^{p+1-i}\pi^{i-1}$. No independence of the
additional $s_p,s_{p+1}$ is asserted or needed.
\end{proof}
% END UF.Conjugation

% BEGIN UF.Scope
\section{Attribution and formal scope}
The rank-three criterion is standard and appears as Proposition 2.1 of
Hou, Feng, Leemans and Qu~\cite{HFLQ}. The inverted-orbit argument
reorganizes their structural proof; the character projectors make explicit
the classical decomposition used in their Lemma 2.4. The cyclotomic
arithmetic is classical. Its use here supplies one proof of the order,
coordinates and automorphisms of the presentation used in their Section 5.
We make no historical priority claim for that organization.

The accompanying application supplies all source-specific proofs. In its
tightness converse it retains the stated odd/even classification of
Cunningham and Pellicer; that input is not a conclusion of this common
paper. Sylow's theorem, Eisenstein's criterion, elementary finite group
theory and the usual determinant norm are standard background. Every
ideal and coordinate calculation used here is proved here.

The unchanged Lean module \texttt{OddKernelIntersection} verifies five
statements: the prime-index intersection lemma and its generation
variant, two reflection identities, and the square of the literal
incorrect marking discussed in the application. It does not formalize
the orbit, character or cyclotomic arguments, or the complete
rank-three criterion. The written proofs of those statements remain
part of this paper's mathematical dependency boundary.
% END UF.Scope
\begin{thebibliography}{9}
\bibitem{HFLQ} D.-D. Hou, Y.-Q. Feng, D. Leemans and H.-P. Qu,
\emph{String C-groups of order $4p^m$}, arXiv:2407.10388v1,
15 July 2024, 12 pp.
\end{thebibliography}
\end{document}
